\section{Physical: operators and read modes}
\label{sec:physical}

\paragraph{Grammar.} A complete standalone program.
\begin{lstlisting}
Physical ::= physical { Decl* Fact* Def* return [canon] r as Head; }
Fact ::= host_set R [discharged];
       | consumer canonical;
Def ::= r = (pipe | let | exists) [canon] Op;
Op  ::= unit | rename R@path {..} [constants {source = literal, ...}]
      | select ... | point select ... | range select ...
      | bind ... | guard ... | project r keep {..} | merge [K] l, r
      | probe [first|smallest] [K] l, r
      | semi [eager|lazy] [K] l, r | anti [eager|lazy] [F] l, r
      | union [any|all] [fused] r1, r2, ... [BlockOrigin]
\end{lstlisting}
Bind retains RA's \code{bind r with [for] Pat [: Type] = (Expr)}
and unchanged Environment suffix: the optional token is semantic, not a read
mode. Generator Bind is streamed unnest, never an independent build target.

\paragraph{Contract.} $\mathsf{PHeading}(t,A)$ checks:
(1)~declarations, unique results, earlier operands, one-to-one renames whose
mapping and literal bindings partition declared columns,
fresh bind columns, subset projections, visible expression uses, a defined and
reachable output;
(2)~probe/merge/semi keys are duplicate-free and equal the operands' common
heading; semi and anti keep the left heading; anti's key is a subset of the
common heading; probe and merge take the union;
(3)~range is a nonempty conjunction of constant/parameter bounds on one Ord column;
(4)~union has at least two operands with equal headings and consistent interface
types; its origin checks free/interface, prefix/domain and witness metadata;
(5)~an all-bare head's dependencies equal the returned heading.
The analysis derives rows, source identifiers, types, computed columns, hosts,
keyed/merge/range reads, scans, modes and canon flags. With
$\mathsf{reads}(r)$ the operand and return uses of $r$:
\begin{mathpar}
\Rule{PH-Pipe}{\mathsf{mode}(r)=pipe\quad r\text{ is not a literal-bearing rename}\\
 \mathsf{reusableHost}(r)\ \lor\
 (\mathsf{linearOrFused}(r)\land r\notin\dom A.key\_reads)}
 {\mathsf{PModeOK}(r)}

\Rule{PH-Let}{\mathsf{mode}(r)\in\{let,exists\}\quad\mathsf{reads}(r)>0}
 {\mathsf{PModeOK}(r)}

\Rule{PH-Block}{\text{branch scopes derived from union origins and fixed operands}\\
 \text{Rust payloads and free-variable binding identities preserved}}
 {\mathsf{blockBindingsOK}(source,target)}

\Rule{WF-Physical}{\mathsf{PHeading}(t,A)\quad\mathsf{PMode}(t,A)}
 {\mathsf{WF}_{Physical}(t)\Downarrow A}
\end{mathpar}
A reusable host is a noncanonical host rename without literal bindings;
literal-bearing renames retain their filtered relation. Linear-or-fused means one read,
or only shared prefix reads inside a checked fused group. \code{exists} needs an
empty heading. \code{probe smallest} needs an adjacent admissible intersection
of $k\ge2$ providers, as specified by \textsc{PH-Smallest}.
Lazy semi/anti needs a reusable host read only lazily on that key. A witness
union (\code{any}) has an empty interface; a fused union names its prefix and
every arm depends on it; cross-arm canon on a statically disjoint union is
rejected. Physical checks neither IX demands nor canonical return evidence.

\paragraph{Smallest candidate intersection.}
Let $p$ be the outer valuation and $v$ one new column. An introducing probe
and its maximal adjacent chain of probe/semijoin membership checks read
providers $R_1,\ldots,R_k$ with headings $K_i\cup\{v\}$, where
$\varnothing\ne K_i\subseteq\mathsf{heading}(p)$ and
$v\notin\mathsf{heading}(p)$. Each provider is retained or a reusable host.
The introducing key is $K_1$; subsequent keys are $K_i\cup\{v\}$.
Every intermediate result is a noncanonical pipe, consumed by the next check.
The final result retains its written mode and canonicalization.
\begin{mathpar}
\Rule{PH-Smallest}{k\ge2\quad C_i(p)=\{v\mid(p[K_i],v)\in R_i\}\\
 j=\mathop{\mathrm{arg\,min}}_{1\le i\le k}|C_i(p)|}
 {\bigcap_{i=1}^{k}C_i(p)
   =\{v\in C_j(p)\mid\forall i\ne j,\ v\in C_i(p)\}}

\Rule{PH-KeyProjection}{J=K\cup\{v\}\subseteq\mathsf{heading}(R)\quad
 R_J=\pi_J(R)\ \text{retained}\\
 R_J\text{ is private to one smallest membership edge}}
 {p\ltimes_J R=p\ltimes_J R_J}
\end{mathpar}
The opt-in \code{smallest_intersections} rewrite records this choice without
altering keys or Rust payloads. Explicit occurrence advice may additionally
retain the exact key projection of a membership provider with extra payload,
by PH-KeyProjection. Later full reads still use the original provider.
The correspondence checker permits only these private projections, used once
by a selected smallest semijoin, and erases them back to the original named
graph. Ordinary empty/\code{all}/result selections do not introduce key views.
Independent interposed relational definitions may move before the entire chain in stable dependency order;
Rust nodes and dependencies on intermediate results block this exposure.
Empty/\code{all}/result-name advice retains its existing selection meaning.
Occurrence advice \code{smallest(R@a,S@b,...)} names every provider of exactly
one complete admissible site, through rename or unary relational provenance.
Repeated, missing, partial, ambiguous or blocked sites are rejected, not widened.
The public Physical contract and rewrite correspondence check the result.
Selection uses only recorded columns and provenance, never Rust type inspection.
Structures supplies a candidate directory and full-key membership access for
every provider; native execution counts each set once and enumerates the minimum
(written-order ties) without allocating a candidate collection per valuation.
In an advised VO intersection, the first Fine subatom introduces $v$ and
subsequent adjacent subatoms check it; they do not independently bind $v$.
E-Check correctly emits semijoins against their full occurrence relations.
PH-KeyProjection supplies candidate-only views for these semijoin sides
without losing payload correlation at later full expansions.


\paragraph{Retained streams.}
The opt-in Physical rewrite \code{cq_rewrites::retain_streams(r,...)} changes
only the mode of named stream results from \code{pipe} to \code{let}; operators,
keys and Rust steps are unchanged.
\begin{mathpar}
\Rule{EXT-RetainStream}{\mathsf{mode}(r)=pipe\quad r\text{ is not a host rename}}
 {\mathsf{mode}(r):=let\text{ preserves erasure and the exhaustive bag of }r}
\end{mathpar}
Under exhaustive pure-total evaluation the retained rows equal those of the
stream, so the rewrite is a refinement checked by \code{check_rewrite}. It is a
decomposition decision: a back end may start a new kernel or task at the
retained rows, so the consumer's parallel domain becomes those rows instead of
the stream's first relation.

\paragraph{Checked host sets.}
The opt-in Physical rewrite \code{cq_rewrites::host_sets(R,...)} consumes
relation-name advice and records \code{host_set R;} in the program text.
Each host-set fact names a used declared host, occurs once, and requires
admissible Leibniz equality on every declared column. It is not a host-type
inference or an unchecked assertion: a certified fact requires the exact
runtime certificate obligation \textsc{IF-HostSet}; a \code{discharged} fact
requires the invoking compiler's proof (\textsc{CTX-Discharge}).
Host-set and consumer facts may be interleaved in the fact prefix.
Section~\ref{sec:context} specifies their admission and discharge.
\begin{mathpar}
\Rule{PH-HostSet}{\code{host\_set R}\in t\quad r=\mathsf{rename}(R)\quad
 H=\mathsf{heading}(r)}
 {\mathsf{key}(r)=H}
\end{mathpar}
Join combines operand keys; projection preserves a key only when it retains
every key column. A returned heading containing the inherited key needs no
\code{canon}; a projection dropping that key still does. The rewrite removes
only boundaries implied by this transfer, leaving all other decisions intact.
The current \code{smallest} intersection law denotes Set. In bag/K-relation
semantics, replacing candidate enumeration by membership also requires this
checked set fact or a derived uniqueness key on the candidate source;
multiplicity gating belongs to the future K-relation contract.

\paragraph{Meaning.} Erasure to RA prime: probe and merge are joins, semi is a
semijoin, point/range are selections, modes and canon flags are dropped, union
and anti keep their origins. Context facts erase after their certificate or
discharge obligations; modes and choices change evaluation, not meaning.
Erasure retains the scalar/generator distinction, including in IX demand and
Physical rewrite certificates. Item multiplicity cannot justify scalar
uniqueness or literal inference.

\paragraph{Recorded decisions.} \code{pipe} (one streamed read), \code{let}
(retained, read more than once), \code{exists} (retained witness); \code{canon}
boundaries; \code{probe first|smallest}; \code{merge}; \code{semi/anti
eager|lazy}; \code{union any|all}, \code{fused} prefixes and dedup placement.
Defaults: written branch order, test unions \code{any}, binding unions
\code{canon} unless proved disjoint, fused prefixes for per-row arms. A fused
group streams its prefix once; fallback uses if-empty control flow and never
evaluates a later Rust arm after a witness.
In the preceding RA a natural join's left operand streams and its right
operand is indexed and probed; the default plan keeps written order, so the
first clause of each scope is the stream. Physical erasure remains exact; it does
not silently commute operands or invent demands.

\paragraph{Rewrite check (\code{cq_lang_physical::check_rewrite}).} Both programs
pass the contract; declarations, head, returned result and the definition graph
of named logical operators and operands are preserved; execution choices and
modes may change; recorded join and predicate placement decisions and
branch-scoped free-variable bindings are kept (PH-Block). It has no IX source and does not establish final
canonicality.
